% Corrigé pédagogique original, fondé sur « Organisation de courses de fond » (6 p.).
% Aucun diagramme du professeur reproduit ; choix de modélisation justifiés.
\section{Cas final : organisation de courses de fond -- correction complète des 9 travaux}
\begin{memo}
Cette partie est la \textbf{référence de travail prioritaire}. Elle répond \textbf{aux neuf demandes} du dernier TD. Le texte du professeur fait autorité. Ce corrigé est \textbf{une proposition argumentée}, pas une correction officielle. Il faut déclarer nos hypothèses H1--H4 quand les règles ne sont pas explicites.
\end{memo}
\textbf{Faits à retenir :} course libre dès 16 ans, semi-marathon dès 18 ans, marathon dès 20 ans ; triathlon et vélo : inscription par club, en ligne ou par courrier, avec licence au lieu de certificat. Préinscription internet jusqu'à J--15, sous réserve de places. Sur place : dossier terminé au moins 1 h avant le départ, pas de triathlon/vélo. Par courrier : terminé au moins J--5. Paiement invalide ou certificat non conforme : annulation en ligne ; par courrier, demande de document et maintien à l'étape adaptée. Club classé par ses cinq meilleurs temps, sinon disqualifié (au plus quatre arrivants). SMS en option payante à 1 euro. Transpondeur : numéro de dossard, temps et passages, restitution obligatoire.

\begin{warningbox}
\textbf{Hypothèses déclarées :} H1 : un agent du secrétariat saisit les inscriptions reçues sur papier ou au guichet ; H2 : on sépare l'état de l'\texttt{Inscription} du support de paiement ; H3 : le statut d'inscription suit des étapes communes, sauf les parcours plus rapides qui sautent plusieurs étapes ; H4 : la catégorie d'âge et l'admissibilité sont vérifiées par \texttt{Epreuve} (le texte n'impose aucune classe responsable). Ne transformez pas ces choix en « faits du sujet ».
\end{warningbox}

\subsection{1. Cas d'utilisation : acteurs, objectifs et paquetages}
\textbf{Acteurs primaires :} Concurrent, Club (pour vélo et triathlon), AgentSecrétariat, ServiceCourses, Administrateur (si gestion technique précisée). \textbf{Acteurs secondaires :} Banque/PrestataireCarte, SystèmeChronométrage (antennes et puces externes, sous H5 si le périmètre du logiciel ne les inclut pas), ServiceSMS (si externe). Un transpondeur passif n'est \textbf{pas automatiquement} acteur ; seul le dispositif qui échange effectivement des messages avec le système peut l'être.
\begin{center}
\begin{tikzpicture}[scale=.84,transform shape,node distance=9mm and 17mm]
\node[actor] (cc) at (0,2.8) {Concurrent};\node[actor] (club) at (0,.8) {Club};
\node[actor] (sec) at (0,-1.3) {Secrétariat};\node[actor] (srv) at (0,-3.3) {Service courses};
\node[usecase] (online) at (4.5,3.2) {S'inscrire sur internet};
\node[usecase] (cour) at (4.5,1.7) {Inscrire par courrier};
\node[usecase] (place) at (4.5,.2) {Inscrire sur place};
\node[usecase] (consult) at (4.5,-1.3) {Rechercher un classement};
\node[usecase] (classer) at (4.5,-2.8) {Enregistrer classement};
\node[usecase] (result) at (4.5,-4.3) {Enregistrer les temps};
\node[actor] (bank) at (8.7,3.4) {Banque};
\node[actor] (chrono) at (8.7,-4.3) {Système de\\chronométrage};
\draw[navy] (2.7,-5.05) rectangle (7.1,4.05);
\node[font=\scriptsize,fill=white] at (4.9,4.04){Système d'organisation des courses};
\draw[navy] (cc)--(online);\draw[navy](cc)--(consult);
\draw[navy](club)--(online);\draw[navy](club)--(cour);
\draw[navy](sec)--(cour);\draw[navy](sec)--(place);
\draw[navy](srv)--(classer);\draw[navy](srv)--(result);
\draw[navy](bank)--(online);\draw[navy](chrono)--(result);
\end{tikzpicture}
\end{center}
\textbf{Paquetages fonctionnels :} \texttt{Inscriptions}, \texttt{EpreuvesEtClubs}, \texttt{Classements}, \texttt{Chronometrage}. Un paquetage de cas n'est pas une étape du scénario. Les cas « consulter conditions » et « télécharger formulaire » appartiennent à l'espace public ; « gérer récompenses » au ServiceCourses. \textbf{Attention :} un Club n'est pas un « Concurrent amélioré » ; c'est un rôle institutionnel différent.

\subsection{2. Description textuelle détaillée : inscription sur internet}
\begin{memo}
\textbf{Méthode exacte d'ASI\_C1, page 36 :} Identification (nom, but, acteurs) ; Séquencement (déclencheur, préconditions, enchaînement nominal, postconditions) ; variantes et exceptions ; éventuellement exigences non fonctionnelles et IHM. Le TD de course de fond demande explicitement ce travail : ne pas le réduire à un diagramme.
\end{memo}
\begin{tabularx}{\linewidth}{@{}p{40mm}X@{}}
\toprule
\textbf{Nom du cas} & S'inscrire à une course sur internet.\\
\textbf{But} & Obtenir une inscription définitive après paiement accepté et vérification des pièces.\\
\textbf{Acteur principal} & Concurrent pour une course individuelle ; Club de la discipline pour triathlon ou vélo.\\
\textbf{Acteurs secondaires} & Banque (carte) ; Secrétariat (chèque, certificat ou licence).\\
\textbf{Déclencheur} & Le concurrent choisit une course et demande son inscription depuis le site web.\\
\textbf{Préconditions} & Une épreuve est publiée et accessible via le site. On \emph{contrôle pendant le scénario} la date, les places et l'admissibilité.\\
\textbf{Postcondition succès} & Une \texttt{Inscription} est \texttt{DEFINITIVE} ; paiement accepté et justificatif conforme.\\
\textbf{Postconditions sans succès} & Avant la préinscription, aucun dossier si quota/date/âge refusé. Après paiement invalide ou pièce non conforme : \texttt{ANNULEE}, concurrent prévenu. Dossier temporairement \texttt{PREINSCRIPTION} ou \texttt{FINANCIERE} en attendant les pièces.\\
\bottomrule
\end{tabularx}
\textbf{Enchaînement nominal (course individuelle ordinaire, carte) :}
\begin{enumerate}[label=\textbf{\arabic*.}]
\item Le concurrent sélectionne la course. Le système affiche ses conditions, sa date et son tarif.
\item Le système vérifie le quota, la limite J--15 et la condition d'âge liée à la discipline.
\item Le concurrent fournit ou confirme ses données personnelles. Le système retrouve les données déjà connues, si elles existent.
\item Le système crée l'inscription au statut \texttt{PREINSCRIPTION}.
\item Le concurrent choisit de payer par carte. Le système transmet la demande de paiement à la banque.
\item La banque accepte. Le système fait passer l'inscription au statut \texttt{FINANCIERE}.
\item Le concurrent envoie par courrier un certificat médical de moins d'un an. Le secrétariat le reçoit et le contrôle.
\item Si le certificat est conforme, le système passe à \texttt{DEFINITIVE} et le concurrent est informé.
\end{enumerate}
\textbf{Enchaînements alternatifs :}
\begin{itemize}
\item \textbf{A1 [étape 3 : concurrent déjà connu] :} le système propose les anciennes informations ; le concurrent les valide ou les corrige ; retour à l'étape 4.
\item \textbf{A2 [étapes 5--6 : chèque] :} le concurrent envoie son paiement par courrier ; le secrétariat le contrôle ; le statut devient \texttt{FINANCIERE} quand le chèque est accepté. Le certificat est vérifié ensuite.
\item \textbf{A3 [triathlon / vélo] :} le \texttt{Club} inscrit son membre par internet ; on contrôle l'adhésion au club spécialisé ; la \texttt{Licence} conforme remplace le certificat médical ; il n'y a \textbf{pas} d'inscription sur place.
\end{itemize}
\textbf{Exceptions :}
\begin{itemize}
\item \textbf{E1 [étape 2 : places insuffisantes ou date J--15 dépassée] :} refus avant préinscription. La condition d'âge empêche également l'inscription aux courses individuelles non autorisées.
\item \textbf{E2 [étape 6 : paiement non valide] :} inscription annulée et concurrent informé.
\item \textbf{E3 [étape 8 : certificat non conforme] :} inscription annulée et concurrent informé.
\end{itemize}
\textbf{Règles métier et remarques :} libre \(\geq16\) ans, semi \(\geq18\), marathon \(\geq20\) ; quota et date J--15 ; certificat médical datant de moins d'un an (ou licence du club spécialisé). \textbf{N'inventez pas} un upload internet de certificat ni une double autorisation/débit : ici la validation définitive est liée à la vérification du justificatif.\par
\textbf{Contrainte non fonctionnelle possible (suggestion, non donnée mot pour mot) :} limiter l'accès aux données personnelles et médicales aux personnes autorisées.
\begin{warningbox}
Un \textbf{paiement accepté} ne signifie pas \textbf{inscription définitive}. La préinscription, le paiement et le contrôle du justificatif sont \emph{trois étapes différentes}. Un courrier incomplet, lui, n'est pas automatiquement annulé : le secrétariat relance le concurrent et conserve l'étape en cours.
\end{warningbox}
\begin{tabularx}{\linewidth}{@{}p{29mm}XXX@{}}
\toprule
Pour comparer & Internet & Sur place & Courrier \\
\midrule
Date clé & Préinscription au plus tard J--15 & Dossier complet 1 h avant départ & Dossier complet au plus tard J--5\\
Qui intervient ? & Concurrent (ou club), banque, secrétariat & Concurrent, secrétariat & Concurrent/club, secrétariat\\
Mode paiement & Carte ou chèque & Espèces, carte ou chèque & Règlement joint, souvent chèque\\
Justificatif & Certificat envoyé par courrier ; licence si club & Certificat remis et contrôlé & Certificat joint ; licence si club\\
Si pièce manque & Dossier attend tant que possible & Validation immédiate impossible & Conserver étape et demander pièce\\
Triathlon/vélo & Club uniquement & \textbf{Non admis} & Club uniquement\\
\bottomrule
\end{tabularx}
\source{Organisation de courses de fond, pages 3--4 ; ASI\_C1, page 36.}

\subsection{3. Les trois séquences système N1 d'inscription}
\textbf{Règle :} un seul participant \texttt{:SystèmeCourses}, plus acteurs externes. Aucun contrôleur, écran ou DAO. Chaque case \texttt{alt} correspond à un véritable embranchement.
\begin{center}
\begin{tikzpicture}[x=.94cm,y=.49cm,scale=.91,transform shape]
\foreach \x/\t in {0/Concurrent,3.5/SystèmeCourses,7.3/Banque,11/Secrétariat}{\node[dbox] at (\x,0){\t};\draw[life](\x,-.5)--(\x,-20.8);}
\draw[assoc](0,-1)--node[above,font=\scriptsize]{1: choisirEpreuve()}(3.5,-1);
\draw[<-,dashed](0,-2)--node[above,font=\scriptsize]{2: conditions, quota, tarif}(3.5,-2);
\draw[teal](-.6,-2.6)rectangle(4.3,-4.6);
\node[font=\tiny,anchor=west] at (-.4,-2.87){break [quota/date/âge non valides]};
\draw[<-,dashed](0,-3.73)--node[above,font=\tiny]{3: refus motivé, pas de dossier}(3.5,-3.73);
\draw[assoc](0,-5.2)--node[above,font=\scriptsize]{4: saisirOuConfirmerDonnees()}(3.5,-5.2);
\draw[<-,dashed](0,-6.2)--node[above,font=\scriptsize]{5: préinscription enregistrée}(3.5,-6.2);
\draw[assoc](0,-7.2)--node[above,font=\scriptsize]{6: choisirPaiement()}(3.5,-7.2);
\draw[teal](-.6,-7.7)rectangle(11.9,-12.25);
\node[font=\tiny,anchor=west] at (-.4,-7.94){alt [carte]};
\draw[assoc](3.5,-8.75)--node[above,font=\scriptsize]{7a: demanderPaiement()}(7.3,-8.75);
\draw[<-,dashed](3.5,-9.7)--node[above,font=\scriptsize]{8a: accepté/refusé}(7.3,-9.7);
\draw[teal,densely dashed](-.6,-10.18)--(11.9,-10.18);
\node[font=\tiny,anchor=west] at (-.4,-10.4){[chèque]};
\draw[assoc](11,-11.48)--node[above,font=\scriptsize]{7b: saisirRésultatContrôleChèque()}(3.5,-11.48);
\draw[teal](-.6,-12.75)rectangle(11.9,-16.2);
\node[font=\tiny,anchor=west] at (-.4,-13){alt [paiement invalide]};
\draw[<-,dashed](0,-13.85)--node[above,font=\tiny]{9a: ANNULEE, avertissement}(3.5,-13.85);
\draw[teal,densely dashed](-.6,-14.35)--(11.9,-14.35);
\node[font=\tiny,anchor=west] at (-.4,-14.58){[paiement accepté]};
\draw[<-,dashed](0,-15.4)--node[above,font=\tiny]{9b: FINANCIERE, attendre pièce}(3.5,-15.4);
\draw[assoc](11,-16.75)--node[above,font=\scriptsize]{10: contrôlerCertificatOuLicence()}(3.5,-16.75);
\draw[teal](-.6,-17.2)rectangle(11.9,-20.45);
\node[font=\tiny,anchor=west] at (-.4,-17.43){alt [justificatif conforme]};
\draw[<-,dashed](0,-18.25)--node[above,font=\tiny]{11a: DEFINITIVE, avertissement}(3.5,-18.25);
\draw[teal,densely dashed](-.6,-18.66)--(11.9,-18.66);
\node[font=\tiny,anchor=west] at (-.4,-18.87){[justificatif non conforme]};
\draw[<-,dashed](0,-19.76)--node[above,font=\tiny]{11b: ANNULEE, avertissement}(3.5,-19.76);
\end{tikzpicture}
\end{center}
\textbf{N1 internet -- lecture exacte :} les deux \texttt{alt} « paiement » et « justificatif » sont deux décisions \textbf{successives}. La seconde n'a lieu que si le paiement a été accepté. Le secrétariat peut saisir l'information après réception du certificat \emph{par courrier}. Le club remplace l'acteur principal et présente une \emph{licence} pour triathlon/vélo. Les réponses aux contrôles font partie du N1, même si les échanges postaux eux-mêmes ont lieu hors système.
\begin{center}\begin{tikzpicture}[x=1cm,y=.72cm]
\foreach \x/\t in {0/Concurrent,5/SystèmeCourses,10/Secrétariat}{\node[dbox] at (\x,0){\t};\draw[life](\x,-.5)--(\x,-8.6);}
\draw[assoc](0,-1)--node[above,font=\scriptsize]{1: remettreFormulaireEtPièces()}(10,-1);
\draw[assoc](10,-2)--node[above,font=\scriptsize]{2: saisirInscription()}(5,-2);
\draw[densely dashed,-Stealth](5,-3)--node[above,font=\scriptsize]{3: résultat contrôles}(10,-3);
\draw[teal] (4,-3.6) rectangle (10.75,-8.2);
\node[font=\scriptsize,anchor=west] at (4.2,-3.95){alt [valide, avant départ moins 1h et places]};
\draw[assoc](10,-4.7)--node[above,font=\scriptsize]{4a: enregistrer définitivement()}(5,-4.7);
\draw[densely dashed,-Stealth](5,-5.7)--node[above,font=\scriptsize]{5a: confirmer inscription}(10,-5.7);
\draw[teal,densely dashed](4,-6.1)--(10.75,-6.1);
\node[font=\scriptsize,anchor=west] at (4.2,-6.4){[sinon]};
\draw[densely dashed,-Stealth](5,-7.25)--node[above,font=\scriptsize]{4b: rejet et motif}(10,-7.25);
\end{tikzpicture}\end{center}
\textbf{N1 sur place :} la branche est refusée aussi pour triathlon/vélo. Pour un paiement carte, la Banque reste un acteur secondaire distinct (ajouter son échange avant confirmation).
\begin{center}\begin{tikzpicture}[x=1cm,y=.72cm]
\foreach \x/\t in {0/Concurrent,5/Secrétariat,10/SystèmeCourses}{\node[dbox] at (\x,0){\t};\draw[life](\x,-.5)--(\x,-8.1);}
\draw[assoc](0,-1)--node[above,font=\scriptsize]{1: envoyerDossierPapier()}(5,-1);
\draw[assoc](5,-2)--node[above,font=\scriptsize]{2: saisirEtContrôlerDossier()}(10,-2);
\draw[densely dashed,-Stealth](10,-3)--node[above,font=\scriptsize]{3: état du dossier}(5,-3);
\draw[teal](4.2,-3.6)rectangle(10.6,-7.8);
\node[font=\scriptsize,anchor=west] at (4.4,-3.9){alt [complet, reçu avant J--5]};
\draw[densely dashed,-Stealth](10,-4.7)--node[above,font=\scriptsize]{4a: inscription définitive}(5,-4.7);
\draw[teal,densely dashed](4.2,-5.15)--(10.6,-5.15);
\node[font=\scriptsize,anchor=west] at (4.4,-5.46){[pièce manquante ou paiement non conforme]};
\draw[densely dashed,-Stealth](10,-6.25)--node[above,font=\scriptsize]{4b: maintien état adéquat}(5,-6.25);
\draw[assoc](5,-7.05)--node[above,font=\scriptsize]{5b: demander pièce au concurrent}(0,-7.05);
\end{tikzpicture}\end{center}
\textbf{N1 courrier :} la règle J--5 s'applique à la validation finale ; une pièce manquante peut être relancée, sans annulation automatique immédiate. Pour une inscription de club, le courrier vient du Club.

\subsection{4. Classes du domaine : noms, propriétés, opérations, relations}
\begin{center}\begin{tikzpicture}[scale=.86,transform shape]
\node[dbox,align=left,minimum width=35mm] (co) at (0,3) {\textbf{Concurrent}\\nom : String\\naissance : Date\\sexe : Sexe\\estMajeurPour(e) : Bool};
\node[dbox,align=left,minimum width=37mm] (ins) at (5,3) {\textbf{Inscription}\\date : Date\\dossard : String\\statut : StatutInscription\\valider() / annuler()};
\node[dbox,align=left,minimum width=40mm] (ep) at (10,3) {\textbf{Epreuve}\\dateDépart : DateTime\\quotaMax : Integer\\tarif : Money\\estAdmissible(c) : Bool};
\node[dbox,align=left,minimum width=35mm] (club) at (0,0) {\textbf{Club}\\nom : String\\discipline : Discipline};
\node[dbox,align=left,minimum width=35mm] (pa) at (5,0) {\textbf{Paiement}\\montant : Money\\mode : ModePaiement\\estAccepté() : Bool};
\node[dbox,align=left,minimum width=37mm] (just) at (10,0) {\textbf{Justificatif}\\dateValidité : Date\\type : TypeJustif\\estConforme() : Bool};
\node[dbox,align=left,minimum width=35mm] (temps) at (0,-3) {\textbf{Passage}\\instant : DateTime\\tempsNet : Durée};
\node[dbox,align=left,minimum width=35mm] (trans) at (5,-3) {\textbf{Transpondeur}\\codeDossard : String\\estRestitué : Bool};
\node[dbox,align=left,minimum width=37mm] (cla) at (10,-3) {\textbf{Classement}\\position : Integer\\tempsFinal : Durée\\calculerPosition()};
\draw[navy](co)--node[above,font=\scriptsize,pos=.10]{1}node[above,font=\scriptsize,pos=.86]{0..*}(ins);
\draw[navy](ins)--node[above,font=\scriptsize,pos=.10]{0..*}node[above,font=\scriptsize,pos=.86]{1}(ep);
\draw[navy](co)--node[left,font=\scriptsize]{0..1 / 0..*}(club);
\draw[navy](ins)--node[left,font=\scriptsize]{1 / 0..1}(pa);
\draw[navy](ins) to[bend left=13] node[right,font=\scriptsize,pos=.5]{1 / 0..1} (just);
\draw[navy](ins)--node[above,sloped,font=\scriptsize]{0..1}(trans);
\draw[navy](trans)--node[above,font=\scriptsize]{0..*}(temps);
\draw[navy](ins)--node[above,sloped,font=\scriptsize]{0..1}(cla);
\end{tikzpicture}\end{center}
\textbf{Précision critique :} les multiplicités abrégées sur les liens obliques ci-dessus se lisent comme des \textbf{hypothèses de modélisation}, à détailler et justifier : une Inscription concerne exactement un Concurrent et une Epreuve ; un Concurrent peut avoir plusieurs Inscription ; un Paiement est rattaché au plus à une inscription individuelle dans ce modèle simplifié. Une catégorie de classement est une \textbf{valeur métier}, pas une nouvelle personne. Un Transpondeur peut être réaffecté dans le temps : les passages restent attachés à l'Inscription (ajouter une association précise en version de production).
\textbf{Classes supplémentaires utiles :} \texttt{Edition}, \texttt{Troncon}, \texttt{Antenne}, \texttt{Recompense}, \texttt{Notification}, \texttt{ClubClassement}. Ne pas stocker un rang figé avant l'arrivée : la position est calculée depuis les temps. \textbf{Classe d'association} : \texttt{Inscription} porte les attributs du lien Concurrent--Epreuve. Les contraintes d'âge et de quotas appartiennent à \texttt{Epreuve} ; les statuts à \texttt{Inscription}.

\textbf{Multiplicités à maîtriser pour ne pas perdre de points :}\par\smallskip
\begin{tabularx}{\linewidth}{@{}p{55mm}p{20mm}X@{}}
\toprule
Association (de A vers B) & A vers B & Lecture et hypothèse \\
\midrule
\texttt{Concurrent -- Inscription} & \texttt{0..* / 1} & Un concurrent peut s'inscrire plusieurs fois ; une inscription correspond à un concurrent.\\
\texttt{Epreuve -- Inscription} & \texttt{0..* / 1} & Une épreuve reçoit plusieurs inscriptions ; chacune vise une seule épreuve.\\
\texttt{Inscription -- Paiement} & \texttt{0..1 / 1} & Un dossier peut attendre le règlement ; un règlement enregistré appartient à un dossier (H2).\\
\texttt{Inscription -- Justificatif} & \texttt{0..* / 1} & Un dossier peut avoir des pièces successives ; chaque pièce est reliée à un dossier (H7).\\
\texttt{Epreuve -- Troncon} & \texttt{1..* / 1} & Parcours divisé en tronçons ; un tronçon est rattaché à son épreuve (H8).\\
\texttt{Troncon -- Antenne} & \texttt{1 / 1} & Une antenne à la fin de chaque tronçon : simplification du TD (H8).\\
\texttt{Transpondeur -- Passage} & \texttt{0..* / 1} & Plusieurs passages peuvent provenir d'une puce ; un passage porte son émetteur.\\
\texttt{Inscription -- Classement} & \texttt{0..1 / 1} & Un classement individuel est lié à une inscription terminée ; aucun rang avant résultat (H9).\\
\bottomrule
\end{tabularx}
\textbf{Sens de la table :} \texttt{A -- B : 0..* / 1} signifie qu'un A peut avoir plusieurs B et qu'un B est rattaché à un seul A. Ce sont des \textbf{choix de modélisation justifiés}, pas des multiplicités toutes écrites dans l'énoncé.\\
\textbf{Ne pas confondre :} \texttt{Inscription} rattache un \texttt{Concurrent} à une \texttt{Epreuve}. Un \texttt{Club} n'est pas un type de \texttt{Concurrent}. Le classement \textbf{par club} est calculé avec les cinq meilleurs temps ; si au plus quatre membres terminent, le club est disqualifié.

\textbf{Vue métier complémentaire -- chronométrage et récompenses :}
\begin{center}
\begin{tikzpicture}[scale=.86,transform shape]
\node[dbox] (ep) at (0,1.6) {Epreuve};
\node[dbox] (tr) at (3,1.6) {Troncon};
\node[dbox] (ant) at (6,1.6) {Antenne};
\node[dbox] (pas) at (9,1.6) {Passage};
\node[dbox] (trans) at (9,-.1) {Transpondeur};
\node[dbox] (ins) at (3,-.1) {Inscription};
\node[dbox] (clas) at (3,-1.8) {Classement};
\node[dbox] (rew) at (7.5,-1.8) {Recompense};
\draw[navy](ep)--node[above,font=\tiny,pos=.25]{1}node[above,font=\tiny,pos=.78]{1..*}(tr);
\draw[navy](tr)--node[above,font=\tiny,pos=.22]{1}node[above,font=\tiny,pos=.80]{1}(ant);
\draw[navy](ant)--node[above,font=\tiny,pos=.22]{1}node[above,font=\tiny,pos=.80]{0..*}(pas);
\draw[navy](trans)--node[right,font=\tiny]{1 / 0..*}(pas);
\draw[navy](ins)--node[left,font=\tiny]{1 / 0..1}(clas);
\draw[navy](clas)--node[above,font=\tiny]{0..* (H9)}(rew);
\end{tikzpicture}
\end{center}
\textbf{Noter H8--H9 :} un transpondeur ne remplace pas la classe \texttt{Inscription}. Pour relier chaque temps au concurrent de la bonne course, une affectation de puce à inscription est nécessaire (association ou classe d'association si l'affectation porte des dates).\par
\subsection{5. Paquetages et dépendances}
\begin{center}\begin{tikzpicture}[node distance=12mm and 20mm]
\node[cbox] (s) {\textbf{Inscriptions}\\Concurrent, Inscription, Paiement\\Justificatif};
\node[cbox,right=of s] (ev) {\textbf{Epreuves et clubs}\\Epreuve, Club, Edition};
\node[cbox,below=of s] (ch) {\textbf{Chronométrage}\\Transpondeur, Passage\\Troncon, Antenne};
\node[cbox,right=of ch] (res) {\textbf{Classements}\\Classement, Récompense};
\draw[umlarrow](s)--(ev);\draw[umlarrow](res)--(s);\draw[umlarrow](res)--(ch);\draw[umlarrow](ch)--(ev);
\end{tikzpicture}\end{center}
\textbf{Critère :} découpage par responsabilité métier ; les dépendances pointent vers le paquetage \emph{utilisé}. Pas de dépendance circulaire. Si une opération impose un cycle, déplacer un concept dans un noyau partagé ou revoir le sens du lien.

\subsection{6. États de l'inscription individuelle avec déclencheurs complets}
\begin{center}\begin{tikzpicture}[node distance=15mm and 20mm]
\node[state] (pre) {Préinscription};
\node[state,right=of pre] (fin) {Inscription financière};
\node[state,right=of fin] (def) {Inscription définitive};
\node[state,below=of fin] (ann) {Annulée};
\node[circle,fill=navy,inner sep=2pt,left=10mm of pre] (ini) {};
\node[draw,circle,minimum size=7mm,inner sep=0pt,right=14mm of def] (end) {}; \node[fill=navy,circle,minimum size=3mm,inner sep=0pt] at (end.center) {};
\draw[assoc](ini)--node[above,font=\tiny]{préinscrire [quota et date OK]}(pre);
\draw[assoc](pre)--node[above,font=\tiny]{paiementAccepté / passerFinancière()}(fin);
\draw[assoc](fin)--node[above,font=\tiny]{pièceConforme / confirmer()}(def);
\draw[assoc](def)--(end);
\draw[assoc](pre) to[bend right=28] node[below,font=\tiny]{paiementRefusé / annuler()} (ann);
\draw[assoc](fin)--node[right,font=\tiny]{certificatRefusé / annuler()}(ann);
\end{tikzpicture}\end{center}
\textbf{Énumération :} \texttt{StatutInscription = \{PREINSCRIPTION, FINANCIERE, DEFINITIVE, ANNULEE\}}. Dans le parcours \textbf{courrier}, si une pièce manque on conserve PREINSCRIPTION ou FINANCIERE et on peut reprendre le contrôle : représenter une \textbf{boucle de régularisation} \texttt{pièceReçue [conforme] / continuer()}. Dans le parcours \textbf{sur place}, trois étapes peuvent être franchies dans la même opération (pas forcément trois écrans). \textbf{Délai :} J--15 bloque l'entrée internet, J--5 la validation courrier et 1 h le guichet ; ne pas inventer un \texttt{after(24h)} inexistant dans ce TD. Le mot \texttt{after(durée)} du cours est pertinent seulement quand un délai relatif déclenche réellement un changement d'état.

\subsection{7. Deux DCP : inscription internet et recherche classement}
\begin{center}\begin{tikzpicture}[node distance=10mm and 12mm]
\node[actor] (a) {Concurrent};
\node[cbox,right=of a] (v) {«dialogue»\\PageInscription};
\node[cbox,right=of v] (c) {«contrôle»\\ControleInscription};
\node[cbox,right=of c,yshift=26mm] (e) {«entité»\\Epreuve};
\node[cbox,right=of c,yshift=13mm] (i) {«entité»\\Inscription};
\node[cbox,right=of c,yshift=0mm] (con) {«entité»\\Concurrent};
\node[cbox,right=of c,yshift=-13mm] (p) {«entité»\\Paiement};
\node[cbox,right=of c,yshift=-26mm] (j) {«entité»\\Justificatif};
\draw[navy](a)--(v);\draw[navy](v)--(c);\draw[navy](c)--(e);\draw[navy](c)--(i);\draw[navy](c)--(con);\draw[navy](c)--(p);\draw[navy](c)--(j);
\end{tikzpicture}\end{center}
\textbf{DCP inscription :} cinq entités du modèle : \texttt{Concurrent}, \texttt{Epreuve}, \texttt{Inscription}, \texttt{Paiement}, \texttt{Justificatif}. \texttt{PagePaiement} et \texttt{PageRecapitulatif} sont des dialogues \textbf{supposés} (H6), non fournis comme maquettes. La seule maquette réellement donnée est celle de recherche des classements. Le contrôle orchestre ; les entités valident la règle métier.
\begin{center}\begin{tikzpicture}[node distance=9mm and 12mm]
\node[actor] (a) {Visiteur};\node[cbox,right=of a] (v) {«dialogue»\\PageRecherche};
\node[cbox,right=of v] (c) {«contrôle»\\ControleRecherche};
\node[cbox,right=of c,yshift=16mm] (e) {«entité»\\Concurrent};
\node[cbox,right=of c] (i) {«entité»\\Inscription};
\node[cbox,right=of c,yshift=-16mm] (cl) {«entité»\\Classement};
\draw[navy](a)--(v);\draw[navy](v)--(c);\draw[navy](c)--(e);\draw[navy](c)--(i);\draw[navy](c)--(cl);
\end{tikzpicture}\end{center}
\textbf{DCP classement :} les critères exacts de la maquette sont \texttt{Dossard, Nom, Prénom, Sexe, Catégorie}. Les entités doivent reprendre exactement les noms du domaine : \texttt{Concurrent, Inscription, Classement}. \textbf{Interdit :} un lien direct Dialogue--Entité. Un \texttt{ControleRecherche} distinct est justifié car c'est un autre cas d'utilisation.

\subsection{8. Communication pour rechercher un classement}
\begin{center}\begin{tikzpicture}[node distance=15mm and 22mm]
\node[actor] (v) {Visiteur};\node[cbox,right=of v] (vue) {:PageRecherche};
\node[cbox,right=of vue] (ctrl) {:ControleRecherche};
\node[cbox,right=of ctrl] (dom) {:Classement};
\draw[navy](v)--node[above,font=\tiny]{1: saisirCritères()}(vue);
\draw[navy](vue)--node[above,font=\tiny]{2: rechercher()}node[below,font=\tiny]{5: afficher()}(ctrl);
\draw[navy](ctrl)--node[above,font=\tiny]{3: obtenirRésultats()}node[below,font=\tiny]{4: résultats}(dom);
\end{tikzpicture}\end{center}
\textbf{Complément de modèle :} \texttt{ControleRecherche} peut demander au \texttt{Concurrent} les identités et à \texttt{Inscription} les dossards. L'ordre des messages est écrit \textbf{sur les liens}, pas sur un axe du temps. Une numérotation détaillée des appels imbriqués \texttt{1, 1.1, 1.2} est aussi valide si cohérente.

\subsection{9. Séquence N3 MVC : inscription sur place}
\begin{center}\begin{tikzpicture}[x=.82cm,y=.54cm]
\foreach \x/\t in {0/Agent,3/PageGuichet,6/ControleInscription,9/Epreuve,12/Inscription}{\node[dbox,minimum width=21mm] at (\x,0){\t};\draw[life](\x,-.5)--(\x,-12.5);}
\draw[assoc](0,-1)--node[above,font=\tiny]{1: saisirDossier()}(3,-1);
\draw[assoc](3,-2)--node[above,font=\tiny]{2: inscrireSurPlace()}(6,-2);
\draw[assoc](6,-3)--node[above,font=\tiny]{3: verifierQuotaAgeDélai()}(9,-3);
\draw[densely dashed,-Stealth](9,-4)--node[above,font=\tiny]{4: résultat}(6,-4);
\draw[teal](5.25,-4.6)rectangle(12.8,-11.9);
\node[font=\tiny,anchor=west] at (5.45,-4.9){alt [admissible]};
\draw[assoc](6,-5.8)--node[above,font=\tiny]{5a: contrôlerPaiementPièce()}(12,-5.8);
\draw[densely dashed,-Stealth](12,-6.75)--node[above,font=\tiny]{6a: statut contrôle}(6,-6.75);
\draw[assoc](6,-7.7)--node[above,font=\tiny]{7a: confirmer()}(12,-7.7);
\draw[assoc](6,-8.6)--node[above,font=\tiny]{8a: afficherConfirmation()}(3,-8.6);
\draw[teal,densely dashed](5.25,-9.2)--(12.8,-9.2);
\node[font=\tiny,anchor=west] at (5.45,-9.5){[quota atteint / délai / justificatif invalide]};
\draw[assoc](6,-10.4)--node[above,font=\tiny]{5b: afficherRefus(motif)}(3,-10.4);
\end{tikzpicture}\end{center}
\textbf{Participants métier à conserver dans une version exhaustive :} \texttt{Concurrent}, \texttt{Paiement}, \texttt{Justificatif}, au besoin \texttt{Banque} pour carte et \texttt{Club} en exclusion du guichet. Le contrôle appelle les \emph{opérations métier des entités}; il ne calcule pas lui-même les seuils d'âge et de délai. L'interface ne parle jamais directement à \texttt{Epreuve} ni à \texttt{Inscription}. Le N3 reprend les interactions externes du N1.

\textbf{N3 détaillée -- messages qui doivent exister, même si le schéma précédent est compact :}
\begin{longtable}{@{}p{13mm}p{52mm}p{80mm}@{}}
\toprule
No & Émetteur vers destinataire & Message et rôle \\
\midrule\endhead
1 & Agent \(\to\) \texttt{PageGuichet} & \texttt{saisirDossier(...)}. Vue : saisie des informations et pièces.\\
2 & Vue \(\to\) \texttt{ControleInscription} & \texttt{inscrireSurPlace(...)} ; le contrôle orchestre.\\
3 & Contrôle \(\to\) \texttt{Concurrent} & \texttt{retrouverOuCreer(...)} ; réutiliser les données d'une précédente course.\\
4 & Contrôle \(\to\) \texttt{Epreuve} & \texttt{verifierAdmissibilite()}. Règles du métier : \textbf{1 h}, quota et âge ; \textbf{pas triathlon/vélo}.\\
5a & \texttt{Epreuve} \(\to\) Contrôle & \texttt{refus(motif)} ; contrôle \(\to\) vue \texttt{afficherRefus(motif)}. Fin du cas si inadmissible.\\
5b & Contrôle \(\to\) \texttt{Paiement} & \texttt{verifier(mode,montant)} ; en cas de carte, solliciter aussi la \texttt{Banque} externe.\\
6b & Contrôle \(\to\) \texttt{Justificatif} & \texttt{estConforme()} pour le certificat médical récent.\\
7b & Contrôle \(\to\) \texttt{Inscription} & \texttt{enregistrerEtValider()} seulement si \textbf{tous les contrôles sont OK}.\\
8b & Contrôle \(\to\) \texttt{PageGuichet} & \texttt{afficherConfirmation(dossard)} puis vue \(\to\) Agent.\\
\bottomrule
\end{longtable}
\textbf{Branches à écrire dans la N3 :} \texttt{alt [épreuve non admissible] / [admissible]} ; puis \texttt{alt [paiement invalide OU certificat invalide] / [tous conformes]}. En cas d'erreur, Le contrôle informe la vue ; le dossier ne devient pas \texttt{DEFINITIVE}. \textbf{H10 :} la manière de tracer une saisie incomplète n'est pas spécifiée dans le texte ; ne pas inventer une notification SMS pour le guichet.\par
\textbf{Cohérence forte :} conserver les mêmes noms de dialogues, contrôles et entités entre le DCP et la N3. Les noms \texttt{Concurrent}, \texttt{Epreuve}, \texttt{Paiement}, \texttt{Justificatif} et \texttt{Inscription} reprennent ceux du domaine.

\subsection{Dernière vérification des neuf réponses}
\begin{tabularx}{\linewidth}{@{}p{24mm}X@{}}\toprule
Questions & Critère de cohérence à vérifier avant de rendre \\\midrule
1--2 & Même acteur principal, même scénario, différences club / individuel claires.\\
3 & Trois canaux modélisés, cas d'erreur et retours, système noir.\\
4--5 & Noms métier et multiplicités ; aucune classe de DAO dans le domaine, pas de cycle.\\
6 & Énumération alignée et branches d'annulation/régularisation présentes.\\
7--8 & Dialogue--Contrôle--Entités, acteurs à l'extérieur, mêmes noms, messages numérotés.\\
9 & N3 MVC détaillée, mêmes événements que N1, contrôle court, règles dans les entités.\\\bottomrule\end{tabularx}